\documentclass{article}
\usepackage{fullpage}
\usepackage{amsmath,amssymb,amsfonts}
\usepackage{graphicx}
\usepackage{natbib}

\def\R{{\mathbb{R}}}
\def\pr{{\rm Pr}}
\def\E{{\mathbb E}}
\def\X{{\mathcal X}}
\def\Y{{\mathcal Y}}
\def\H{{\mathcal H}}
\def\G{{\mathcal G}}
\def\B{{\mathcal B}}
\def\bias{{\rm bias}}
\def\supp{{\rm supp}}

\newtheorem{thm}{Theorem}
\newtheorem{lemma}[thm]{Lemma}
\newtheorem{cor}[thm]{Corollary}
\newtheorem{claim}[thm]{Claim}
\newtheorem{defn}[thm]{Definition}
\newtheorem{assump}{Assumption}
\newtheorem{open}{Open problem}
\newenvironment{proof}{\noindent {\sc Proof:}}{$\Box$ \medskip}

\DeclareMathOperator*{\argmax}{arg\,max}

\title{Active learning using the near-Neighborhoods algorithm}

\begin{document}

\maketitle

\section{Setup}

Basic Transductive setup:
\begin{itemize}
\item Intance set $\X =\{x_1,\dots,x_n\}$, $x_i \in {\bf X}$, where
  ${\bf X}$ is a metric space. . The set
  $\X$ is fixed and known to the learner. The probability of each $x
  \in \R$ is $1/n$.
\item Label space $\Y = \{-1,+1\}$. Each $x \in \X$ is associated with
  a conditional probability function $\eta(x) = \pr(Y=1|X=x)$.
\item Let $g^*$ denote the Bayes-optimal classifier on $\X$:
  $$ g^*(x) = \mbox{sign}(2 \eta(x) - 1) .$$
  \item Label Queries: the algorithm can select any example $x \in
    \X$, as a response it gets a label in $\Y$ chosen acccording to
    $\eta(x)$.    
  \end{itemize}

\section{Notation}
\begin{itemize}
\item {\bf Balls:} For any $x \in \X$ an $r\geq0$,
  $B(x,r)$ denotes the subset of $\X$ inside ball of radius $r$ centered at $x$.
  \item {\bf Probability of a ball} The probability of a ball is the
    fraction of points in $\X$ that fall inside $B(x,r)$.
    $$ \mu(B(x,r)) = \frac{ |B(x,r)|}{n}$$
  \item {\bf conditional probability given a ball:}
    $$ \eta(B(x,r)) = \frac{\sum_{x \in B(x,r)}\eta(z)}{|B(x,r)|}$$
  \item {\bf measuring radius in terms of probability}
For any $x \in X$ and any $p > 0$, let $r_p(x)$ denote the smallest
radius $r$ such that $\mu(B(x, r)) \geq p$.
\end{itemize}

\newcommand{\NNB}{{\bf NNB}}
\newcommand{\ANNB}{{\bf ANNB}}
\newcommand{\ONNB}{{\bf oracle-NNB}}

\section{The near-neighborhoods algorithm}
The near neighborhood (\NNB) algorithm is a variant of AKNN where the balls
that are considered for the prediction of a point $x$ are {\em all} balls
the contain $x$ while in AKNN only balls that are centered at $x$ are
considered.

Suppose that we are given $n$ labeled examples chosen uniformly at
random from $\X$ and labeled according to $\eta(x)$.

Given that set of labels each ball is as $-,?,+$ according to the
condition used in AKNN. We say that a ball labeled $\pm$ is {\em
  determined}, a ball labeled $?$ is {\em undetermined}.

We define the partial order on balls that is defined by containment.

The labels of the balls induce a labeling on the points in $\X$. Each
point has one of {\em four} labels: $-1,+1,?,!$. Given a point $x \in
\X$ we define the {\em minimal set of determined balls for $x$},
denoted $\B(x)$, to be the set of determined balls that contain $x$
that is minimal (in terms of the containment partial order). We can
construct this set by starting with the set of all determined balls containing
$x$ and eliminating any ball that has a subset-ball that is also in
the set.

Given the set of balls $\B(x)$ the label associated with $x$
according to the following rules:
\begin{itemize}
\item If $\B(x)$ is empty the label is undetermined, or $?$.
\item If all of the balls in $\B(x)$ are labeled $+$, the label is $+1$.
\item If all of the balls in $\B(x)$ are labeled $-$, the label is $-1$
\item If $\B(x)$ contains both labels, then the label is {\em
    over-determined} or $!$. We will also refer to over-determined
  examples as the ``known unknown''.
\end{itemize}

We define the numerical prediction $\hat{y}$ to be $0$  if the label
is undetermined or overdetermined. We set it equal to the label otherwise

We define the error of the prediction as the produce
$-\hat{y}{y}$. Note that the error can be $-1,0,+1$. 

\NNB\ is a passive learning algorithm, in other words the queriy
points are selected IID from $\X$.

\begin{claim}
  The algorithm \NNB\ converges to Bayes rule as the number of
  labeled examples increases to $\infty$.
  \end{claim}

\section{Active \NNB}

The main idea behind active \NNB\, or \ANNB\, is to concentrate the
queries on the known unknown. However, this is not done to the
exclusion of the other parts of the space, as any example can have an
unbounded number of mind changes. In other words, an example whose label at some point is $+1$
can be labeled $-1$ later on and vice versa. (These are the ``unknown
unknowns'').

\ANNB\ operates in phases. At the start of phase $t$ a concentration
area $C_t \subset \X$ is computed. Then, $2n$ queries are are
made. There are $n$ {\em passive} queries and $n$ {\em active} queries.
The passive queries are selected uniformly at random from $\X$, the
active quries are selected uniformly at random from $C_t$.

After the answers to the queries are recieved, they are used to update
the counter $k_+,k_-$ in the balls. Passive points are used in each
ball that contains them. Active points are a bit more subtle, they are
used only in balls that are conained in $C_t$. In other words if a
ball contains an element $x$ that is {\em not} in $C_t$ then this ball
cannot use {\em any} of the active samples for epoch $t$.

Based on the additional queries, the label $\pm +1,?$ for each ball is
updated, and based on that the label of each $x \in \X$ is updated.

Finally, $C_{t+1}$ is calculated as follows. Let $U_t$ be the set $x
\in \X$ such that the label of $x$ is overdetermined, or ``!''.  $C_{t+1}$
is defined as the union of the minimal sets of determined balls for $x
\in U_t$, In other words:
\[
  C_{t+1} = \bigcup_{x \in U_t} \B(x)
  \]

  \section{A comparator algorithm}
  In order to bound the query complexity of \ANNB\ we consider a
  variant of \ANNB\ that is aided by and oracle. Intuitively, the
  oracle removes the need for a passive sample by informing the
  algorithm that the prediction on a point $x$ is ``final''. I.e. it
  will not change as the number of queries increases to infinity.

  For this algorithm, which we denote by \ONNB\, we can more easily
  give bounds. As this framework is local, our goal is to provide
  point-wise bounds.

  \ONNB\ differs from \ANNB\ in two important ways
  \begin{enumerate}
  \item It makes only active queries, not passive queries.
  \item The set $U_t$ contains, in addition to overdetermined
    points. The points where the label is $+1$ or $-1$ but the oracle
    informs the algorithm that there will be mind cahnges on $x$ in
    the future. 
  \end{enumerate}

  Intuitively, \ONNB\ contines to collect information on $x$ until the
  label of $x$ converges, at which point it does not need more
  samples for $\B(x)$ so it stop collecting this information.

  \section{Pointwise analysis}

  Fix a point $x \in \X$ and compare the behaviour of \NNB\ and \ONNB\
  on $x$.

  To make life easier, we consider the {\em expected number of}
  queries made on each point. (actually, need to disregard queries
  that are part of a batch that does not contain the ball). Denote by.
    
  
\end{document}
